% SPDX-License-Identifier: GPL-3.0-or-later
% Copyright (C) 2026 Alessandro Greco
% Русское издание: Unicode/LuaLaTeX и русская типографика.
\usepackage{fontspec}
\usepackage{polyglossia}
\setdefaultlanguage{russian}
\setotherlanguage{english}
\usepackage[autostyle=once]{csquotes}
\setmainfont{Libertinus Serif}
\setsansfont{Libertinus Sans}[Scale=MatchLowercase]
\setmonofont{DejaVu Sans Mono}[Scale=0.85]
% Моноширинный шрифт с кириллицей и для внутренних переключений шрифтов tex4ht (EPUB):
% без него polyglossia сообщает об ошибке для ec-lmtt, а символы теряются. Без переносов,
% как у \setmonofont: иначе пути вроде labs/ru/лаборатория разбивались бы дефисом.
\newfontfamily\cyrillicfonttt{DejaVu Sans Mono}[Scale=0.85,HyphenChar=None]
\newfontfamily\cyrillicfontsf{Libertinus Sans}[Scale=MatchLowercase]
% В EPUB tex4ht местами работает в кодировке OT1, где нет кавычек «» и „:
% задаём их по умолчанию как символы Unicode, чтобы они не терялись.
\ifdefined\HCode
  \ProvideTextCommandDefault{\guillemotleft}{«}
  \ProvideTextCommandDefault{\guillemotright}{»}
  \ProvideTextCommandDefault{\quotedblbase}{„}
\fi
\usepackage[a4paper,inner=3cm,outer=2.5cm,top=2.5cm,bottom=3cm]{geometry}
\usepackage{microtype}
\setlength{\emergencystretch}{3em}
\usepackage{xcolor}
\usepackage{graphicx}
\usepackage{booktabs,tabularx,longtable,multirow}
\usepackage{amsmath,amssymb}
\usepackage{siunitx}
% Десятичная запятая и тонкий пробел: несуществующий профиль locale=RU не нужен.
\sisetup{output-decimal-marker={,},group-separator={\,}}
\usepackage{enumitem}
\usepackage{tikz}
\usetikzlibrary{arrows.meta,positioning,shapes,calc,fit,backgrounds,decorations.pathreplacing,decorations.pathmorphing,matrix,chains}
\usepackage{bytefield}
\usepackage{listings}
% listings печатает «-» и «*» как моноширинные символы, только если текущее семейство
% совпадает с \ttdefault; в русском тексте polyglossia берёт \cyrillicfonttt, и в
% листингах (PDF и EPUB) появлялись математические «−» и «∗». Все листинги книги
% набираются моноширинным шрифтом, поэтому всегда выбираем моноширинный вариант.
\makeatletter
% \relax — как в самом listings: иначе «\char95» и следующие за ним цифры («_332»)
% читаются как один номер символа.
\def\lst@ttfamily#1#2{#1\relax}
\makeatother
\usepackage[most]{tcolorbox}
\usepackage[backend=biber,style=alphabetic,maxbibnames=99]{biblatex}
\addbibresource{../shared/bibliography/tor.bib}
% Общая библиография написана для итальянского издания: здесь заменяются примечания на
% итальянском. «Consultato il ...» превращается в url + urldate, которые biblatex выводит по-русски.
\DeclareSourcemap{\maps[datatype=bibtex,overwrite=true]{
  \map{
    \step[fieldsource=note, match={Consultato il 23 settembre 2026}, final]
    \step[fieldsource=howpublished, fieldtarget=url]
    \step[fieldsource=url, match=\regexp{\A\\url\{(.+)\}\z}, replace={$1}]
    \step[fieldset=urldate, fieldvalue={2026-09-23}]
    \step[fieldset=note, null]
  }
  \map{\step[fieldsource=entrykey, match=\regexp{\Atorblogwebtunnel\z}, final]\step[fieldset=note, fieldvalue={12 марта 2024 г.}]}
  \map{\step[fieldsource=entrykey, match=\regexp{\Aperry2018torbrowser\z}, final]\step[fieldset=note, fieldvalue={Черновик}]}
  \map{\step[fieldsource=entrykey, match=\regexp{\Adsa2022\z}, final]\step[fieldset=note, fieldvalue={Официальный журнал ЕС L 277 от 27.10.2022}]}
  \map{\step[fieldsource=entrykey, match=\regexp{\Acodiceprivacy\z}, final]\step[fieldset=note, fieldvalue={Сводный текст, опубликованный итальянским органом по защите персональных данных (Garante per la protezione dei dati personali)}]}
  \map{\step[fieldsource=entrykey, match=\regexp{\Awinter2014spoiled\z}, final]\step[fieldset=note, fieldvalue={позднее вошла, с участием других авторов, в \emph{Privacy Enhancing Technologies Symposium} (PETS 2014)}]}
  \map{\step[fieldsource=entrykey, match=\regexp{\Alegge167\z}, final]\step[fieldset=note, fieldvalue={Gazzetta Ufficiale, Serie Generale №~277 от 27.11.2017}]}
  \map{\step[fieldsource=entrykey, match=\regexp{\Anistpqc2024\z}, final]\step[fieldset=note, fieldvalue={три постквантовых стандарта, опубликованные одновременно}]}
  \map{\step[fieldsource=entrykey, match=\regexp{\Ahayden2014metadata\z}, final]\step[fieldset=note, fieldvalue={Johns Hopkins Foreign Affairs Symposium, 1 апреля 2014 г.}]}
  \map{\step[fieldsource=entrykey, match=\regexp{\Apfitzmann2010terminology\z}, final]\step[fieldset=howpublished, fieldvalue={Версия 0.34}]}
}}
\DeclareBibliographyAlias{legislation}{misc}
\usepackage{imakeidx}
\makeindex[intoc,title={Предметный указатель}]
\usepackage[hidelinks]{hyperref}
\usepackage[nameinlink]{cleveref}
% cleveref не склоняет названия: получалось «в глава 7», «из раздел 8». Внутри фразы
% используются несклоняемые сокращения («в гл. 7», «из разд. 8», «в прил. A»), как
% «рис.» и «табл.»; в начале предложения (\Cref) — полные формы в именительном падеже.
\crefname{chapter}{гл.}{гл.}
\Crefname{chapter}{Глава}{Главы}
\crefname{section}{разд.}{разд.}
\Crefname{section}{Раздел}{Разделы}
\crefname{figure}{рис.}{рис.}
\Crefname{figure}{Рис.}{Рис.}
\crefname{table}{табл.}{табл.}
\Crefname{table}{Табл.}{Табл.}
\crefname{appendice}{прил.}{прил.}
\Crefname{appendice}{Приложение}{Приложения}
% Союзы в перечнях ссылок пишем символами Unicode: в EPUB tex4ht местами работает
% в кодировке OT1, где команда \cyri недоступна, и «и» между ссылками терялось.
\newcommand{\crefpairconjunction}{ и }
\newcommand{\crefmiddleconjunction}{, }
\newcommand{\creflastconjunction}{ и }
\newcommand{\crefpairgroupconjunction}{ и }
\newcommand{\crefmiddlegroupconjunction}{, }
\newcommand{\creflastgroupconjunction}{ и }
